\subsection{Interpretation of functions as distributions}

PROPOSITION 6. --- Let ν be a measure on U. The restriction of ν to D(U) is a distribution, which is zero if and only if the measure ν is zero.

Let K be a compact subset of U. For every function \(\varphi\in\mathcal{C}_{\mathrm{K}}^{\infty}(\mathrm{U})\) , we have \(|\langle\nu,\varphi\rangle|\leqslant p_{0,\mathrm{K}}(\varphi)|\nu|(\mathrm{K})\) , hence the restriction of \(\nu\) to \(\mathcal{D}(\mathrm{U})\) is continuous. Since \(\mathcal{D}(\mathrm{U})\) is dense in \(\mathcal{K}(\mathrm{U})\) , by proposition 4(a) of IV, p. 202, the restriction of \(\nu\) to \(\mathcal{D}(\mathrm{U})\) is zero if and only if \(\nu\) is zero.

We will identify the space \(\mathcal{M}(\mathrm{U};\mathbf{C})\) of complex measures on U with a subspace of \(\mathcal{D}'(\mathrm{U})\) . We will also identify the space \(\mathrm{L}_{\mathrm{loc}}^{1}(\mathrm{U})\) with a subspace of \(\mathcal{D}'(\mathrm{U})\) by the map induced by passing to quotients of the map that assigns to \(f\in \mathcal{L}_{\mathrm{loc}}^{1}(\mathrm{U})\) the measure \(f\cdot \mu\) (INT, V, p. 44, § 5, no. 2, def. 2) . In other words, for \(f\in \mathrm{L}_{\mathrm{loc}}^{1}(\mathrm{U})\) and \(\varphi \in \mathcal{D}(\mathrm{U})\) , we have

\[
\langle f, \varphi \rangle = \int_ {\mathrm{U}} f \varphi d \mu .
\]

In particular, for \(p \in [1, +\infty]\) this allows one to identify the space \(\mathrm{L}^p(\mathrm{U})\) with a subspace of \(\mathcal{D}'(\mathrm{U})\) (cf. INT, V, p. 43, § 5, no. 1).

PROPOSITION 7. --- Let \(p \in [1, +\infty]\) . The injection of \(\mathrm{L}^p(\mathrm{U})\) into \(\mathcal{D}'(\mathrm{U})\) is continuous.

Let \(f \in \mathrm{L}^{p}(\mathrm{U})\) . Let K be a compact subset of U, and denote by \(\varphi_{K}\) its characteristic function. For every test function \(\varphi\) with support contained in K, Hölder's inequality implies

\[
| \langle f, \varphi \rangle | = \left| \int_ {\mathrm{U}} f \varphi d \mu \right| \leqslant \mathrm{N} _ {p} (f) \mathrm{N} _ {q} (\varphi) \leqslant \mathrm{N} _ {p} (f) \mathrm{N} _ {q} (\varphi_ {\mathrm{K}}) p _ {0, \mathrm{K}} (\varphi)
\]

where q is the conjugate exponent of p.

In particular, one can identify \(\mathcal{D}(\mathrm{U})\) with a subspace of \(\mathcal{D}'(\mathrm{U})\) .

PROPOSITION 8. --- The space \(\mathcal{D}(\mathrm{U})\) is dense in \(\mathcal{D}'(\mathrm{U})\) .

Let \(\lambda\) be a linear form on \(\mathcal{D}'(\mathrm{U})\) vanishing on \(\mathcal{D}(\mathrm{U})\) . Since \(\mathcal{D}(\mathrm{U})\) is reflexive, there exists a test function \(\varphi \in \mathcal{D}(\mathrm{U})\) such that \(\lambda(f) = \langle f, \varphi \rangle\) for all \(f \in \mathcal{D}'(\mathrm{U})\) . We obtain

\[
0 = \lambda (\overline {{\varphi}}) = \langle \overline {{\varphi}}, \varphi \rangle = \int_ {\mathrm{U}} | \varphi | ^ {2} d \mu ,
\]

whence \(\varphi = 0\) . The proposition then follows from the Hahn--Banach theorem (EVT, II, p. 49, cor. 3 (ii)).

For \(h \in \mathcal{C}^{\infty}(\mathrm{U})\) , the transpose of the continuous linear map \(\varphi \mapsto h\varphi\) of \(\mathcal{D}(\mathrm{U})\) into itself is a continuous linear map from \(\mathcal{D}'(\mathrm{U})\) into itself, which is denoted \(f \mapsto hf\) . This definition is justified because if \(f\) is the distribution associated with a measure \(\nu\) on \(\mathrm{U}\) , then \(hf\) is associated with the measure \(h \cdot \nu\) . Indeed, for every test function \(\varphi \in \mathcal{D}(\mathrm{U})\) , one computes

\[
\langle h f, \varphi \rangle = \langle f, h \varphi \rangle = \int_ {\mathrm{U}} h \varphi d \nu = \int_ {\mathrm{U}} \varphi d (h \cdot \nu).
\]

